\chapter[Project Management]{Project Management}{Project Management}\label{projm}


\rule{0.49\textwidth}{0.75pt} $_{\bigcirc}$ \rule{0.49\textwidth}{0.75pt}\\





\section{Feasibility Study}

\subsection{Technical Feasibility}
TrustChain has a high technical feasibility with well established and widely adopted technologies used in the frontend and backend layers. The frontend will be constructed in HTML5, CSS3 and React.js which will provide a responsive and easy to use interface to register, access and manage the dashboard. It can be integrated with Web3.js or Ethers.js to allow the user interface to interact with blockchain smart contracts easily. This will provide a smooth communication and real time updates of identity verification workflows within the system.

On the backend, there is the use of Node.js and Express.js that deal with the server side operations, API management and communication between the various parts of the system. IPFS integration (Pinata) is a way to achieve the decentralization of identity documents storage which enhances their availability and fault tolerance. Decentralized identifiers, access control and verification logic are implemented in Solidity based smart contracts running on the blockchain layer with Ethereum. Metamask and WalletConnect are embedded tools to facilitate secure wallet authentication and to interact with the blockchain network effectively.

The system also includes cutting edge cryptographic and identity systems like Self Sovereign Identity and ZK STARK based verification. These technologies provide the identity validation to be done without exposing sensitive information and retain the integrity of proofs with the help of KECCAK 256 hashing. This combination is what results in the system being technically sound and in line with the current standards of security.

Design wise, TrustChain adheres to the human centered design principles to make it user-friendly and easily accessible. The user flow remains minimal and has a straightforward onboarding, upload and access document management. This makes it possible to have even the non technical users to interact with the system effectively without having to jeopardize security or functionality of the system. The interface is to be as user friendly as possible and yet easy to understand every step of the process. This enhances the general usability and promotes usage by a broader group of users.
\subsection{Operational Feasibility}

The operational feasibility of TrustChain revolves around how it is able to operate in real world identity verification situations. The system will act in tandem with the other pre-existing systems like the government systems and telecommunication systems. It will act as a verification layer of middle ground to allow smooth adoption and not necessitate a total overhaul of existing systems and workflows. It minimizes reliance on manual verification processes within the departments. This aids in quick processing and delivery of services. The system can be incorporated without impairing the current workflows.

The interaction process is easy and arranged to ensure users can register, upload documents and access them without being complicated. Features such as guardian based recovery and controlled access ensure reliability in case of user errors or access issues. The system reduces repetition of data entry that enhances efficiency and minimizes overhead costs of operation by the users and requesting entities. The users can simply know and manage their data sharing preferences. This enhances confidence and minimizes misunderstandings when verifying. The platform provides uniform experience to various uses.

Institutionally, TrustChain lessens the burden of managing sensitive identities information through offering proof based verification. Raw documents are not sent to organizations but only validation results making it easier to comply and less risky. Smart contracts provide automation of access control and keep transparent a log of all activities on the blockchain. It also minimizes the chances of data duplication and inconsistencies. The automation of logs enhances operations traceability and auditability. This enhances accountability in organizations that use the system.

The system is also scalable in terms of operations, taking advantage of the use of decentralized storage and blockchain infrastructure. It provides uniform performance without any centralized bottlenecks. This renders TrustChain to be applicable in the implementation in various industries that need secure and effective identity checking. Its decentralized structure allows high availability when it is being used to the fullest. It mitigates risks that are caused by centralization system failures. This ensures uninterrupted access and reliable verification services.
\subsection{Economic Feasibility}
Decentralized infrastructure and use of open source technologies can support the economic viability of TrustChain. It eliminates the cost of having costly centralized storage and data management systems. This reduces first development cost and guarantees scalability and durability.

The system lowers the cost of operation by decreasing the number of times of identity verification and the number of times of handling sensitive documents manually. Organizations do not have to invest much in secure storage infrastructure because data is stored in IPFS. Verification and access control are automated by the use of blockchain smart contracts, cutting down on administrative overhead. Moreover, integration with the existing systems will not incur the expense of replacement of the entire infrastructure and the solution is thus cost effective to adopt.

In the long term, TrustChain will help to minimize financial losses related to identity fraud and data breaches. It reduces the risks that identity data can be misused by making sure that the verification is secure and privacy preserving. The scalable architecture is also economically viable as it minimizes maintenance costs in the long run and can be deployed in large scalable manner across sectors.


\section{Project Management}
\subsection{Timeline Chart}
Below are the key milestones for the TrustChain project, along with an estimated timeline based on standard software development cycles and the methodologies described in our research paper.

\begin{table}[h!]
\centering
\begin{tabular}{|p{4cm}|p{7cm}|p{4cm}|}
\hline
\textbf{Milestone} & \textbf{Description} & \textbf{Estimated Completion Date} \\
\hline
Project Initiation \& Planning & Define project scope, create team, collect requirements & August, 2025 \\
\hline
System Design \& Architecture & Design system layers, identity flow and identity access control & September, 2025 \\
\hline
Frontend \& Backend Development & Develop user interface, backend APIs and integrate wallet authentication & November, 2025 \\
\hline
IPFS \& Hashing Implementation & Deploy a multi layer hashing and decentralized storage & December, 2025 \\
\hline
Blockchain \& Contract Development & Write smart contracts to check identity and access control & January, 2026 \\
\hline
Testing \& Validation & Conduct functional, security and performance testing of the system & February, 2026 \\
\hline
User Feedback \& Iteration & Collect feedback and improve usability and system flow & March, 2026 \\
\hline
Deployment \& Integration & Deploy system and test integration to real world use cases & March, 2026 \\
\hline
Final Review \& Documentation & Complete documentation and finalize project report & April, 2026 \\
\hline
Conference Presentation & Write research paper and present at ICTIS conference. & April, 2026 \\
\hline
\end{tabular}
\caption{TrustChain Project Milestones and Timeline}
\end{table}

\\ 


\section{Project Resources}

\subsection{Hardware Requirements}
TrustChain hardware requirements are specified by the necessity to develop, deploy and ensure identity verification in case of different user roles like developers, users and requesting entities. These requirements guarantee the best performance, system responsiveness and safe use of the platform.

1. Development and Hosting Server Requirements
Processor: Intel Core i7 (or equivalent AMD Ryzen 7) – 8th Gen or higher

RAM: Minimum 16 GB (Recommended: 32 GB for for handling blockchain and cryptographic operations)

Storage: 1 TB SSD (Solid State Drive) for faster data access and efficient IPFS integration

GPU (Optional): Not mandatory, as no heavy machine learning training is involved

Internet: High speed broadband with upload/download speeds of at least 100 Mbps

Power Backup: UPS or generator backup to ensure continuous uptime during outages

2. Client Side (Users and Requesters)
Processor: Intel Core i3 or higher

RAM: Minimum 4 GB (Recommended: 8 GB for smoother performance)

Storage: 250 GB HDD or SSD

Compatibility with Devices: compatible with laptops, smartphones, tablets and desktops.

Display: Minimum resolution of 1366 x 768 pixels

Internet: Stable internet connection with 10 Mbps speed for seamless interaction and verification processes



\subsection{Software Requirements}

1. Operating System\\
Development Environment:
Windows 10/11 (best used when setting up blockchain nodes), macOS (use cross platform testing)
Client Environment (User Side):
Windows 10/11, Android 8.0+, iOS 13+, or any other operating system with a current web browser.

2. Frontend Technologies\\
Language: HTML5, CSS3, JavaScript
Framework: React.js (v18) – to create responsive and dynamic user interfaces.
Styling Libraries: Tailwind CSS or Bootstrap (to have a consistent UI design.)

3. Backend Technologies\\
Runtime Environment: Node.js (v16 or above) to scale on server side.
Framework: Express.js (v4.17.1) to route HTTP requests and API.
Decentralized Storage: IPFS with Pinata integration for storing identity documents.

4. Blockchain Layer\\
Platform: Ethereum Network ( OP Sepolia Testnet to deploy and test)
Smart Contract Language: Solidity to verify identity and access control logic.
Wallet Integration: Metamask or WalletConnect to sign a transaction and authenticate.

5. Identity Layer \& Cryptography\\
Framework: Self Sovereign Identity to decentralized identifiers and verifiable credentials.
Mechanism Proof: ZK STARK privacy preserving identity verification.
Hashing Algorithm: KECCAK 256 to verify the integrity of data and validate proofs.

6. Development Deployment Tools\\
Version Control: Git with GitHub integration.
Package Managers: npm to deal with project dependencies.
IDE or Text Editors: Visual Studio Code or any development environment of choice.



\section{Project Estimation}

\subsection{COCOMO Estimation Model}
In order to determine the development effort of TrustChain, the Basic COCOMO Model is used under Semi Detached category since the project contains moderate complexity with the integration of blockchain, cryptographic operations and a team with mixed levels.
\\
\\
Estimated Size: 14,000 lines of code (14 KLOC)\\
Effort (E): $E = 3.0 \times (14)^{1.12} \approx 57 \text{ person-months}$\\
Development Time (D): $D = 2.5 \times (57)^{0.35} \approx 9 \text{ months}$\\
Team Size (P): $P = \frac{57}{9} \approx 6 \text{ members}$

\begin{table}[h!]
\centering
\begin{tabular}{|l|l|}
\hline
\textbf{Parameter}        & \textbf{Value}       \\
\hline
LOC                      & 14,000             \\
Effort                   & \textasciitilde57 PM    \\
Development Time         & \textasciitilde9 months \\
Team Size                & \textasciitilde6 members \\
\hline
\end{tabular}
\caption{Project Parameters}
\label{tab:project_parameters}
\end{table}


\subsection{Function Point Analysis}
Function Point Analysis (FPA) estimates the size and effort of TrustChain by evaluating its functional components such as user inputs, outputs, storage and blockchain interactions. The method does not depend on the programming language and is oriented towards the functionality of the system.

1. Component Breakdown

\begin{table}[h!]
\centering
\begin{tabular}{|l|c|l|c|c|}
\hline
\textbf{Component}                                   & \textbf{Count} & \textbf{Complexity} & \textbf{Weight} & \textbf{Total FP} \\
\hline
External Inputs (registration, upload, access)              & 6     & Medium     & 4      & 24       \\
External Outputs (verification result, proof)          & 4     & Medium     & 5      & 20       \\
User Inquiries (status check, access logs)             & 3     & Simple     & 3      & 9       \\
Internal Logical Files (user data, hashes) & 4     & Medium     & 7      & 28       \\
External Interface Files (IPFS, APIs)             & 3     & Medium     & 7      & 21       \\
\hline
\textbf{Unadjusted Function Points}              &       &            &        & \textbf{102}  \\
\hline
\end{tabular}
\caption{Function Points Calculation}
\label{tab:function_points}
\end{table}

2. Complexity Adjustment
Based on the system features like security, distributed processing, cryptographic computation and integration complexity, a Complexity Adjustment Factor (CAF) of 1.08 is used.
\\
Adjusted Function Points (AFP): $AFP = 102 \times 1.08 \approx 110$\\

3. Effort Estimation
With an average of 10 hours per function point:

Total Effort = 110 × 10 = 1,100 hours\\

A full-time (90 hours/month/person) 3-person team:

Time: $T = \frac{1100}{3 \times 90} \approx 4 \text{ months}$

\rule{0.49\textwidth}{0.75pt} $_{\bigcirc}$ \rule{0.49\textwidth}{0.75pt}\\
